\section{Related Literature}
\label{sec:literature}

This paper connects four literatures, but its contribution does not rest on claiming novelty for any one of their ingredients.

First, a large environmental-federalism literature studies how decentralized jurisdictions combine fiscal and environmental policies. \citet{OatesSchwab1988} show that local governments can choose capital taxes or subsidies together with environmental quality when competing for industry. \citet{BayindirUpmann1998Emissions} studies interjurisdictional emission-tax competition when local firms compete imperfectly, while \citet{MarkusenMoreyOlewiler1993} and \citet{MarkusenMoreyOlewiler1995} emphasize how environmental policy, imperfect competition, and plant-location incentives interact. These papers establish that decentralization can alter environmental-policy incentives and, importantly for our equilibrium repair, that policy changes can alter market structure. We therefore do not claim that strategic environmental federalism, policy-composition distortions, or boundary-regime analysis are new.

Second, the fiscal-competition literature has long recognized that governments compete with both firm-facing fiscal instruments and productive public inputs. \citet{BayindirUpmann1998PublicGoods} compares tax and public-expenditure competition with industrial public goods. \citet{BenassyQuereGobalrajaTrannoy2007} analyze the trade-off between taxation and public inputs, and \citet{HauptmeierMittermaierRincke2012} study strategic interactions in business taxes and productive public inputs. Their framework is especially close to our reaction-function architecture: neighboring fiscal policies can induce responses in both local taxes and public inputs. Indeed, their theory and empirical results allow a neighbor's tax change to induce both an own-tax adjustment and an infrastructure response. Thus neither cross-instrument reaction functions nor the qualitative pattern ``respond to a rival fiscal change through public inputs'' is new. The distinction we seek is narrower: in our model product substitutability itself moves the same cross-instrument best response through zero at a unique threshold on a certified global-SPNE region.

A recent contribution by \citet{MoritaOkoshi2025} is even closer in policy language. They combine fiscal competition for foreign direct investment with public transport infrastructure and show that fiscal competition can raise infrastructure investment under particular transport-cost conditions. Their model contains subsidies, public infrastructure, product-market competition, and threshold effects. The distinction in the present paper is not the coexistence of those ingredients. It is the narrower theorem that product substitutability reverses a simultaneous cross-policy best response through the firm's investment transmission system, together with a matched same-primitives benchmark in which removing conventional investment eliminates that reversal.

Third, the place-based-policy literature documents the prevalence and consequences of firm-specific incentives. \citet{SlatteryZidar2020} describe state and local business incentives and show that large discretionary subsidies are an important instrument of regional economic policy. This evidence motivates treating targeted firm support as a meaningful policy margin, but our model keeps plant locations fixed in order to isolate the product-market and investment channels rather than bidding for discrete location.

Fourth, the environmental industrial-organization literature studies firms' innovation and abatement decisions under market competition and regulation. Most relevant, \citet{StrandholmEspinolaMunoz2025} allow firms to invest simultaneously in conventional cost-reducing R\&D and green R\&D and show that the two investment margins interact with environmental regulation. Hence dual investment is also not, by itself, our contribution.

The paper's surviving contribution lies at the intersection of these literatures. Multi-instrument local competition is known, dual firm investment is known, and environmental policy under oligopoly is known. In the repaired model, however, product-market substitutability uniquely changes the sign of a government's infrastructure response to a rival's firm-specific green subsidy over a nonempty global-SPNE region. At the canonical primitives, the matched benchmark without conventional competitive investment does not reproduce the reversal. This last statement is deliberately not a universal necessity claim: other environmental primitives can generate switching without conventional investment. The contribution is therefore a narrow theorem about the strategic-feedback network of the combined game, not a claim that instrument switching itself is new.